\cabstract{
参照第一次实验作业要求\cite{lecture01}后，在此报告中完成相应的课后作业。

本报告基于课程第一次大作业要求，使用仅依赖 NumPy 的自实现神经网络库完成多个基础模块与模型的实现与实验验证。
工作包括：实现前向/反向传播的基础层（仿射层、激活函数、卷积与池化等）、归一化方法（BatchNorm/LayerNorm/Spatial BatchNorm）、若干便捷的层组合（affine-relu、conv-relu-pool 等）以及优化器与训练框架（Solver）。
在此基础上实现并训练了若干模型，包括两层全连接网络、三层卷积网络与 DeepConvNet，并在 CIFAR-10 数据集上进行了消融与超参实验。实验结果表明所实现的 DeepConvNet 在验证集上达到约 73.9\% 的验证准确率，模型在提升表达能力的同时存在一定的过拟合倾向，建议通过更强正则化、数据增强与合理的学习率调度来进一步提高泛化性能。

}
\ckeywords{卷积神经网络，批量归一化，Dropout，反向传播，CIFAR-10}
\ekeywords{}
